\uuid{86TM}
\exo7id{3775}
\titre{exo7 3775}
\auteur{quercia}
\organisation{exo7}
\datecreate{2010-03-11}
\isIndication{false}
\isCorrection{false}
\chapitre{Espace euclidien, espace normé}
\sousChapitre{Endomorphismes auto-adjoints}
\cpge{mpsi-mp2i.applications-lineaires.b, pcsi.applications-lineaires.b, ptsi.applications-lineaires.b, mpsi-mp2i.matrices-applications-lineaires.b, pcsi.matrices-applications-lineaires.b, ptsi.matrices-applications-lineaires.b, mp-mpi.endomorphismes-euclidien.f, pc.endomorphismes-euclidiens.d, psi.endomorphismes-euclidien.f, pt.isometries-euclidien.e, mp-mpi.reduction-endomorphismes.b, pc.reduction.a, psi.reduction-endomorphismes.a, pt.reduction-endomorphismes.a, bcpst2.produit-scalaire.c, bcpst2.valeurs-propres-vecteurs-propres.a}
\module{Algèbre}
\niveau{L2}
\difficulte{}

\contenu{
\texte{
Soit $M \in \mathcal{M}_n(\R)$ antisymétrique et $f$ l'endomorphisme de $\R^n$
canoniquement associé à $M$.
}
\begin{enumerate}
    \item \question{Montrer que les valeurs propres de $M$ sont imaginaires pures.}
    \item \question{Montrer que $\mathrm{Ker} f \perp \Im f$. En déduire que $g = f_{|\Im f}$ est un isomorphisme de $\Im f$.}
    \item \question{Montrer que $g^2$ est diagonalisable. En déduire que $\mathrm{rg}(M)$ est pair.}
\end{enumerate}
}
